% Author-facing replacement for the proofs of Claims 6.15 and 6.16.
% Current labels: claim:R3_blocks and claim:R3_Vprime.
% Do not apply until the full actual-host integration theorem is CI green.
% No external Gardiner or rank-three classification is used.

% Claim 6.15: replace only the proof body by:
\begin{proof}
Since every automorphism of $G$ extends to an automorphism of~$H$
preserving $V(G)$, the intersections of the imprimitivity blocks
with $V(G)$ form a block system for $\Aut(G)$. The rook action is
primitive. Indeed, the stabiliser of a vertex has two orbits of
size four on the remaining vertices. A nontrivial block containing
that vertex would have size dividing~$9$, hence size~$3$, which
cannot be a union of the stabiliser orbits.
Thus the intersections are either singletons or all of $V(G)$.
In the first case at least nine host blocks are needed, and since
$|H|=18$ and every block has size at least two, all have size~$2$.
In the second case one block contains all nine selected vertices;
since the partition is nontrivial, all blocks have size~$9$.
\end{proof}

% Claim 6.16: replace only the proof body by:
\begin{proof}
If the blocks have size~$9$, the selected graph fills one block,
and $V'$ is the other. Vertex-transitivity gives an automorphism
taking a vertex of $G$ into $V'$. Invariance of the partition then
maps the whole first block onto the second, yielding
$H[V']\cong G$.

If the blocks have size~$2$, each consists of a selected vertex
and its unique outside mate. Every automorphism of $G$ extends to
an automorphism of $H$ preserving these pairs, so it acts on the
mates by the same permutation. Consequently, $H[V']$ is invariant
under the natural $\Aut(\gridth)$-action. This action has two
orbits on unordered pairs of distinct vertices: the edges and
non-edges of $\gridth$. Therefore $H[V']$ is empty, complete,
$\gridth$, or its complement. Since $H$ is regular and
$|V'|=|V(G)|$, counting degrees over the two halves gives
$|E(H[V'])|=|E(G)|=18$; the empty and complete cases are excluded.
Finally, $\gridth$ is self-complementary, as witnessed by
$(i,j)\mapsto(i+j,i-j)$ over $\mathbb F_3$.
\end{proof}
